\begin{table}[ht]
  \centering
  \small
  % tabularx per mandare a capo i test.
  % Cardinalità centrata (c) e allineamento a bandiera (raggedright) per evitare spazi dilatati
  \begin{tabularx}{\textwidth}{@{\hspace{0.5em}} l c >{\raggedright\arraybackslash}X @{\hspace{0.5em}}}
    \toprule
    \textbf{Fase DAG}                        & \textbf{Cardinalità} & \textbf{Test}                                                                                                                                                                                                                                                                                           \\
    \midrule
    Phase~A (nessuna dipendenza)             & 16 test              &
    \texttt{0.1} $\cdot$ \texttt{0.2} $\cdot$ \texttt{0.3} $\cdot$ \texttt{1.1} $\cdot$ \texttt{1.5} $\cdot$ \texttt{1.6} $\cdot$ \texttt{3.3} $\cdot$ \texttt{4.1} $\cdot$ \texttt{4.2} $\cdot$ \texttt{4.3} $\cdot$ \texttt{6.2} $\cdot$ \texttt{6.4} $\cdot$ \texttt{7.2} $\cdot$ \texttt{ext.0.1.nuclei} $\cdot$ \texttt{ext.1.5.testssl} $\cdot$ \texttt{ext.1.5.sslyze} \\
    \addlinespace
    Phase~B (\texttt{depends\_on = ["1.1"]}) & 2 test               &
    \texttt{1.4} $\cdot$ \texttt{2.1}                                                                                                                                                                                                                                                                                                                                         \\
    \bottomrule
  \end{tabularx}
  \caption{Topologia del \gls{DAG} nell'assessment di validazione.}
  \label{tab:dag-topology}
\end{table}